% Part: history
% Chapter: biographies
% Section: kurt-goedel

\documentclass[../../../include/open-logic-section]{subfiles}

\begin{document}

\olfileid{his}{bio}{god}

\olsection{কুর্ট গ্যোডেল}

\olphoto{goedel-kurt}{কুর্ট গ্যোডেল}

কুর্ট গ্যোডেলের (উচ্চারণ: \textsc{গার}-ড্‌ল) জন্ম
১৯০৬ সালের ২৮ এপ্রিল অস্ট্রো-হাঙ্গেরীয় সাম্রাজ্যের
ব্রুন শহরে (বর্তমানে চেক প্রজাতন্ত্রের ব্রনো)। ছোট্ট
কুর্টেল ছিলেন অনুসন্ধিৎসু ও বুদ্ধিমান; তাই বাড়ির
লোকজন তাঁকে প্রায়ই “Der kleine Herr Warum”
(ছোট্ট প্রশ্নবাবু) বলে ডাকতেন। প্রাথমিক বিদ্যালয়
থেকেই পড়াশোনায় তিনি উৎকর্ষ দেখান; কেবল গণিতেই
তিনি সর্বোচ্চ নম্বরের কম পেয়েছিলেন। দুর্বল স্বাস্থ্যের
কারণে গ্যোডেলকে প্রায়ই বিদ্যালয়ে অনুপস্থিত থাকতে
হতো এবং শরীরচর্চার ক্লাস থেকে তাঁকে অব্যাহতি
দেওয়া হয়েছিল। শৈশবে তাঁর বাতজ্বর ধরা পড়ে।
চিকিৎসকদের মূল্যায়ন ভিন্ন হলেও সারা জীবন তিনি
বিশ্বাস করতেন যে এতে তাঁর হৃদ্‌যন্ত্র স্থায়ীভাবে
ক্ষতিগ্রস্ত হয়েছে।

১৯২৪ সালে গ্যোডেল ভিয়েনা বিশ্ববিদ্যালয়ে পড়া শুরু
করেন এবং ১৯৩০ সালে ডক্টরেট লাভ করেন।
প্রথমে তিনি পদার্থবিদ্যা পড়তে চেয়েছিলেন; কিন্তু
দার্শনিক রুডলফ কার্নাপের প্রভাবেও তাঁর আগ্রহ দ্রুত
গণিত, বিশেষত যুক্তিবিদ্যার দিকে যায়। হান্স হানের
তত্ত্বাবধানে লেখা তাঁর অভিসন্দর্ভে অভিন্নতাসহ
প্রথম-ক্রমের বিধেয় যুক্তিবিদ্যার পূর্ণতা উপপাদ্য
প্রমাণিত হয় \citep{Godel1929}। ১৯৩০ সালে তিনি
তাঁর সবচেয়ে বিখ্যাত ফলগুলি পান—প্রথম ও দ্বিতীয়
অসম্পূর্ণতা উপপাদ্য (প্রকাশিত \citealt{Godel1931}-এ)। ভিয়েনায় থাকাকালে
গ্যোডেল ভিয়েনা সার্কলের সঙ্গে গভীরভাবে যুক্ত ছিলেন।
বিজ্ঞানমনস্ক দার্শনিকদের এই গোষ্ঠীতে কার্নাপও
ছিলেন; গ্যোডেলের ফলগুলি কার্নাপের কাজে বিশেষ
প্রভাব ফেলেছিল।
% BN-SRC-845 TARGET-CORRECTION-BEGIN
\par\noindent\textbf{পাঠ-সংশোধন।} ১৯২৯ অভিসন্দর্ভের সাল; ডক্টরেট পাওয়ার সাল ১৯৩০।
\href{https://www.ias.edu/scholars/godel}{প্রতিষ্ঠানের শিক্ষাগত নথি}
দ্বিতীয় সালটি নিশ্চিত করে। ফল পাওয়ার ১৯৩০ ও প্রকাশের
১৯৩১ সাল আলাদা রাখা হয়েছে।
% BN-SRC-845 TARGET-CORRECTION-END

১৯৩৮ সালে গ্যোডেল অ্যাডেলে নিমবুরস্কিকে বিয়ে
করেন। তাঁর বাবা-মা এতে খুশি হননি: অ্যাডেলে শুধু
তাঁর চেয়ে ছয় বছরের বড় ও পূর্বে বিবাহবিচ্ছিন্নই
ছিলেন না, একটি নৈশক্লাবে নৃত্যশিল্পী হিসেবেও কাজ
করতেন। কিন্তু সামাজিক চাপ গ্যোডেলকে প্রভাবিত
করেনি; তাঁর মৃত্যু পর্যন্ত তাঁদের সুখী দাম্পত্য বজায়
ছিল।

১৯৩৮ সালে নাৎসি জার্মানি অস্ট্রিয়া দখল করার পর
গ্যোডেল ও অ্যাডেলে যুক্তরাষ্ট্রে চলে যান। সেখানে
নিউ জার্সির প্রিন্সটনে ইনস্টিটিউট ফর অ্যাডভান্সড
স্টাডিতে তিনি একটি পদ গ্রহণ করেন। অন্তর্মুখী
ও খানিকটা ব্যতিক্রমী স্বভাবের হলেও প্রিন্সটনে তাঁর
সময়টি ছিল সহযোগিতা ও সৃষ্টিশীল কাজে ভরা।
তিনি সেটতত্ত্ব, দর্শন ও পদার্থবিদ্যার উপর প্রবন্ধ
প্রকাশ করেন। বিশেষত, একই প্রতিষ্ঠানের সহকর্মী
আলবার্ট আইনস্টাইনের সঙ্গে তাঁর গভীর বন্ধুত্ব গড়ে
ওঠে।
% BN-SRC-846: অতিরিক্ত ব্যাখ্যা/অনিশ্চয়তা review-ledger-এ; reader-এ প্রয়োগ নয়।

জীবনের শেষ দিকে গ্যোডেলের মানসিক স্বাস্থ্যের
অবনতি ঘটে। ১৯৭৭ সালে তাঁর স্ত্রী হাসপাতালে
ভর্তি হলে তিনি আর গ্যোডেলের জন্য রান্না করতে
পারতেন না। সারা জীবন মানসিক সমস্যায় ভোগা
গ্যোডেল ক্রমে তীব্র সন্দেহবাতিকের শিকার হন।
বিষ খাওয়ানো হবে—এই ভয়ে তিনি খাবার খেতে
অস্বীকার করেন। ১৯৭৮ সালের ১৪ জানুয়ারি
প্রিন্সটনে অনাহারে তাঁর মৃত্যু হয়।

\begin{reading}
গ্যোডেলের পূর্ণাঙ্গ জীবনী জানতে \citet{Dawson1997}
দেখো। আরও জীবনীমূলক লেখা এবং যুক্তিবিদ্যা ও
দর্শনে গ্যোডেলের অবদান নিয়ে প্রবন্ধের জন্য দেখো
\citet{Wang1990}, \citet{Baaz2011},
\citet{Takeuti2003} ও \citet{Sigmund2007}।

গ্যোডেলের ডক্টরেট অভিসন্দর্ভটি মূল জার্মান ভাষায়
পাওয়া যায় \citep{Godel1929}। অসম্পূর্ণতা
উপপাদ্যগুলি নিয়ে মূল প্রবন্ধটি হলো
\citep{Godel1931}। গ্যোডেলের প্রকাশিত রচনাগুলি
মূল ভাষা ও ইংরেজি অনুবাদে তাঁর \emph{Collected Works}
(সংকলিত রচনাবলি)-এর প্রথম দুই খণ্ড
\citet{Godel1986,Godel1990}-তে পাওয়া যায়।
% BN-SRC-847 TARGET-CORRECTION-BEGIN
\par\noindent\textbf{পাঠ-সংশোধন।} উদ্ধৃত প্রথম দুই খণ্ডে প্রকাশিত রচনা আছে;
সব অপ্রকাশিত রচনাও ওই দুই খণ্ডে আছে বলা সংশোধিত।
দেখো \href{https://academic.oup.com/book/54958/chapter-abstract/422765176}{সম্পাদকদের ভূমিকা}।
% BN-SRC-847 TARGET-CORRECTION-END

গ্যোডেলের অসম্পূর্ণতা উপপাদ্যগুলির বিশদ আলোচনা
জানতে \citet{Smith2013} দেখো। উপপাদ্যগুলি নিয়ে
অনানুষ্ঠানিক দার্শনিক আলোচনা আছে মার্ক
লিনসেনমায়ারের পডকাস্টে
\citep{Linsenmayer2014}।
\end{reading}

\end{document}
